\chapter{Introduction}


HOOK INTRO!!!



Nubia’s geographical proximity to the territory of Ancient Egypt, both occupying the Nile Valley region, allowed the two cultures to have different kinds of relations along the millennia, ranging from trade networks and peaceful coexistence to conquest periods. During the New Kingdom (ca. 1550-1069 BCE), Egypt, then a unified and strong state, controlled and colonized Nubia [p.~369-370]{smithNubianExperienceEgyptian2020}\parencites[p.~157]{törökTwoWorldsFrontier2009}[p.~78]{edwardsNubianArchaeologySudan2004}. After the Egyptian withdraw, denominated the Napatan period (ca. 750–270 BCE), after the name of the kingdom’s capital, the relations between Nubia and Egypt were reversed by the conquest and rule of Kushite kings in Upper Egypt during the 25th Dynasty in the Third Intermediate Period \parencites[p.~285-287]{buzonEntanglementFormationAncient2016}[p.~84]{smithRevengeKushitesAssimilation2013}.


The presence of Egyptian-style objects in Nubian funerary contexts and use of Egyptian cultural markers by the Kushite kings was long misunderstood by early Egyptologists, which viewed Nubia as a periphery reliant on Egypt \parencites[p.~191]{smithBackwaterPuritansRacism2022}. These scholars interpreted Nubian engagement with Egyptian culture as a blind assimilation into a “superior” society \parencites[p.~369]{smithNubianExperienceEgyptian2020}. Even during the period of Kushite domination over Upper Egypt, scholars characterized Napatan culture as merely “a provincial imitation of earlier Egypt” \parencites[p.~292]{wilsonBurdenEgyptInterpretation1951}, and “probably of Theban origin” \parencites[p.~22]{breastedHistoryAncientEgyptians1908}.


However, recent scholarship and technological advances in the field of archaeology have revealed a more complex reality. Research now demonstrates that, since the New Kingdom, Egyptian migrants and Nubians developed a multicultural community, blending local and foreign customs and traditions, that continued to evolve throughout the Napatan period \parencites[p.~286]{buzonStrontiumIsotope87Sr2013}. Crucially, the absence of power relations between the two cultures during the Napatan period―with the exception of almost a hundred years of Kushite reign―means that the adoption of Egyptian-style objects by the Nubians cannot be explained as “acculturation”―the assimilation of a weaker population by a dominant state \parencites[p.~528]{smithNubiaEgyptInteraction1998}. This context is fruitful for research into cultural interactions and the consumption of foreign material without the implicating factor of unequal power dynamics \parencites[~p.18-19]{howleyPowerRelationsAdoption2018}{balandaProtectingMummyReinterpretation2014}. Furthermore, Nubia was part of a wider trade network, with evidence of contact in southern Africa, the Middle East and the Mediterranean world \parencites[p.~208]{smithBackwaterPuritansRacism2022}.


Consequently, new theoretical frameworks have been applied to explain the presence of foreign artifacts in Napatan burials, especially those of royalty. Researchers have employed concepts such as cultural entanglements, prestige goods economy, and elite emulation and legitimization to explain how and why the Kushites selectively incorporated Egyptian elements into their culture \parencites{smithNubiaEgyptInteraction1998, smithRevengeKushitesAssimilation2013, yellinEgyptianReligionIts1995, howleyRoyalPyramidTombs2015, howleyPowerRelationsAdoption2018, balandaDevelopmentRoyalFunerary2020}. Cultural entanglement \parencites{dietlerConsumption2010} has been the most employed theoretical approach, especially for the New Kingdom and Napatan periods \parencites{buzonEntanglementFormationAncient2016, gibbonMorphometricAssessmentAppendicular2016, schraderEverydayLifeCollapse2017, spencerIntroductionHistoryHistoriography2017, retzmannNewKingdomPopulation2019, viethNewKingdomTemple2019, budkaAcrossBorders2Living2020, smithNubianExperienceEgyptian2020, smithBackwaterPuritansRacism2022, buzonTumuliTombosInnovation2023, gasperiniAmaraWestCemeteries2023, buzonExploringIntersectionalIdentities2024}. This framework was first proposed for Nubia by \textcite{vanpeltRevisingEgyptoNubianRelations2013}, based on material culture and iconography from the New Kingdom. Cultural entanglements address questions of selective adoption, adaptation and the intermixing of Nubian and Egyptian traditions, by focusing on the agency of the local population that actively selected foreign objects and gave it new meanings and social functions in the new context \parencites{smithDesertRiverConsumption2014, smithColonialEntanglementsEgyptianization2014, buzonEntanglementFormationAncient2016, lohwasserInterkulturalitätVermittlungVerschränkung2017}.


A significant critique of the cultural entanglements model is its tendency to overlook the internal social diversity of ancient Nubia \parencites{lemosForeignObjectsLocal2020}. The Napatan period exemplifies this complexity, spanning approximately five hundred years and a geographical expanse from Lower Nubia to Meroe, beyond the Fifth Cataract. This extensive temporal and regional scope encompasses significant historical developments, including the aftermath of the Egyptian withdrawal form Lower Nubia, the rise of the Napatan dynasty, the hundred-year reign of the 25th Dynasty Kushite kings, subsequent periods of trade and conflict with Assyrian Egypt, and ongoing religious and political relations with the cult of Amun in Thebes. These variations―social, chronological and regional―have yet to be fully addresses, creating a gap in our understanding of Nubian society during this period.


In order to understand the role of Egyptian-style objects in Napatan Nubia, it is therefore essential to account for all these variations, examining all social groups across different times and spaces. A particular fruitful theoretical framework for this task is the concept of “objectscapes”. First developed as a critique of Romanization, an objectscape is defined as “the material and stylistic properties of a repertoire of objects in a given period and geographic range” \parencites[p.368]{pittsObjectscapesManifestoInvestigating2021}. This approach allows for a more nuanced interpretation of material culture within its specific contexts, moving beyond simplistic models of cultural influence to explore how objects materialize social dynamics.


Amulets, pendants bearing culturally significant motifs, provide an ideal case study for applying this approach. As one of the most prominent grave goods found in Nubian burials, they appear across all social groups and throughout the entire temporal and regional span of the Napatan period. The Egyptian influence on these objects is evident in their choice of motifs and material composition, and their abundance allows for a systematic research of their role within specific burial contexts. By examining the different consumption patterns of these Egyptian-style objects between social groups, we can gain critical insights into the broader social interactions. Variations in access and consumption―shaped by factors such as social position, time period and distance to the Egyptian territory―likely influenced how different groups adopted, adapted and rejected certain elements of amulets. Thus, these choices offer an effective lens though which to analyze social and cultural interactions in Napatan Nubia.


\section{History of research}


Nubia has been the center of studies in very recent history compared to Egypt. Although monuments in Lower Nubia were explored since the 19th century, this was mainly due to the interest in Ancient Egyptian antiquities \parencites[p.~3]{raueIntroduction2019}. Archaeological missions started to be carried out in Sudanese territory from the beginning of the 20th century and the Sudan Antiquities Service, later renamed the National Corporation for Antiquities and Museums (NCAM), was founded in 1905.


The term Nubia has been used to describe different ideas. When used geographically, it refers to the region of modern Sudan and includes the Lower Nubia region, now part of the modern state of Egypt. “Nubian” can describe languages currently used in the Middle Nile and Aswan, as well as groups of people in academic literature, although no record exists of self denomination for this particular term \parencites{williamsNubiaBriefIntroduction2020}. Other applied terminologies are Kush and Kushite, which refer to the Kingdoms of Kush, their cultures and people, commonly understood to encompass the Kingdom of Kerma, and the Napatan and Meroitic Periods. In this dissertation, the terms Nubia and Nubian will be employed to refer to the geographical location—from Lower Nubia to the south of Sudan—, and the people living in this region at that particular period. I will also use Kush, Kushite and Napatan to describe the Napatan Period and its characteristics.


The construction and raising of dams in Aswan led to several salvage and excavation projects \parencites{ahmedHistoryArchaeologicalWork2020}. Between 1907 and 1911, different campaigns were conducted in order to excavate and survey sites from different period of Nubian history, including campaigns carried out by George A. Reisner and Cyril M. Firth, as well as the University of Pennsylvania \parencites{reisnerArchaeologicalSurveyNubia1909, reisnerArchaeologicalSurveyNubia1910, randall-maciverAreika1909, woolleyKaranògRomanoNubian1910, woolleyKaranògTown1911, firthArchaeologicalReportNubia1912, firthArchaeologicalSurveyNubia1915, firthArchaeologicalSurveyNubia1927, junkerBerichtÜberGrabungen1920, junkerErmenneBerichtÜber1925, junkerToschkeBerichtÜber1926}. Work continued in the following years. Between 1916 and 1923, Reisner excavated six royal cemeteries in Lower Nubia, while Francis Llewellyn Griffith directed diggings in the region of Sanam, Faras and Kawa \parencites{griffithOxfordExcavationsNubia1922, griffithOxfordExcavationsNubia1923, griffithOxfordExcavationsNubia1926, macadamTemplesKawa1949}. Another planned heightening of the Aswan Dam led to other campaigns in Lower Nubia between 1929 and 1934, when more cemeteries and fortresses were excavated \parencites{emeryExcavationsSurveyWadi1935, emeryRoyalTombsBallana1938, steindorffAnibaMissionArchéologique1935}.


From 1959, the largest international salvage campaign was conducted with the help of UNESCO, as the construction of the High Adam at Aswan was planned to flood a sizable region of Lower Nubia and, consequently, a large number of archaeological sites. More than thirteen countries participated in the campaign which not only conducted excavations but also relocated monumental structures to unaffected sites, with the temple at Abu Simbel the most famous example \parencites[p.~7-8]{raueIntroduction2019}[p.~10-11]{ahmedHistoryArchaeologicalWork2020}. Since then, several international and Sudanese archaeological campaigns have carried out work in Sudan, enriching our understanding of the country’s rich history. Particularly, the use of developing technology, such as isotope and osteological analysis, have shed light into migration and identity \parencites{buzonStrontiumIsotope87Sr2013, buzonTombosNapatanPeriod2014, schraderIlluminatingNubianDark2014, schraderIntraregional87Sr86Sr2019, gibbonMorphometricAssessmentAppendicular2016}. As a result of these campaigns, many publications have been dedicated to Nubian populations and cultures across its history \parencites{adamsNubiaCorridorAfrica1977, oconnorAncientNubiaEgypts1994, welsbyKingdomKushNapatan1996, morkotBlackPharaohsEgypts2000, smithWretchedKushEthnic2003, edwardsNubianArchaeologySudan2004, raueHandbookAncientNubia2019, emberlingOxfordHandbookAncient2020, breyerNapataUndMeroë2022}.


Particularly, our knowledge of the Napatan period has been greatly benefited from these campaigns. Since Reisner and Griffith excavations, further Napatan burial sites, especially non-elite ones, were found and published \parencites[p.~628-630]{lohwasserNapatanNecropoleisBurial2019}. Although comprehensive research into many of these sites is limited, in the last decades, studies have been carried out in order to re-examine Egyptian-style material culture, revealing new insights and interpretations of the period \parencites{lohwasserAspekteNapatanischenGesellschaft2012, dollIdentificationSignificanceTexts1982, dollRoyalMortuaryCult2014, dollKushiteQueenIrtieru2022, balandaProtectingMummyReinterpretation2014, balandaNapatanFuneraryAmulets2015, balandaDevelopmentRoyalFunerary2020, petacchiLocalAspectsCult2014, petacchiBookDeadNapatan2018, petacchiNapatanTombDecorations2022, howleyRoyalPyramidTombs2015,howleyPowerRelationsAdoption2018, strongReexaminingTombQueen2021}.


\section{What are amulets?}


Amulets are the material focus of this dissertation, therefore, it is important to understand what these objects are and their importance, as it is through them that I aim to examine the cultural and social dynamics in Napatan Nubia. By analyzing their archaeological context, iconography and material composition, amulets can be recognized as a materialization of cultural and social structures.


Amulets are one of the most numerous grave goods found in the Nile Valley, especially Egypt, and have been the focus studies, especially those on religion and burial customs. Egyptologists offer differing definitions of these objects, reflecting varied perspectives of their nature. While some emphasize function, others focus on their material characteristics. \textcites[p.~6]{petrieAmulets1914} defined amulets broadly as “any object with means of suspension on the person, and without immediate use or ornamentation”. He also proposed a functional typology followed by other researchers \parencites{andrewsAmuletsAncientEgypt1994}, cf. \parencites[p.~17]{müller-winklerÄgyptischenObjektAmuletteMit1987}. \textcite[p.~20]{müller-winklerÄgyptischenObjektAmuletteMit1987}. defined an amulet as "an object that is supposed to protect the wearer with its magical powers to keep them healthy", while \textcites[p.~2]{herrmannÄgyptischenAmuletteSammlungen2003} rovided a more functional definition as "a small object intended to protect the wearer with magical powers and shield them from evil, provide them with health and other goods, and also allow them to share its magical power." \textcites[p.~6]{andrewsAmuletsAncientEgypt1994}. concurred with the protective function, stating that amulets are “personal ornaments that, because of their shape, color and material, are believed to confer on the wearer, by magical means, certain powers and abilities.” \textcites[p.~19]{dubielAnthropomorpheAmuletteGrabern2004} focused on physical characteristics, defining amulets as small objects with holes or suspension mechanisms that represent a living being or an object. Recently, \textcite[p.~1-4]{quackAltägyptischeAmuletteUnd2022} characterized amulets as wearable objects which the owner believes will give them protection by “supernatural powers”. Quack bases his definition on the denominations the Egyptians themselves gave to these objects, such as “protection”, “to be safe”,  and “to protect”.


While written Egyptian sources and the names given to amulets strongly suggest their protective function, this interpretation cannot be universally applied. Literacy in Ancient Egypt was restricted to the elite and, therefore, these texts were not consumed by the wider population. Furthermore, as \textcites{régenWhenBookDead2010, régenArchaeologyBookDead2017} demonstrated with the case of magical bricks in funerary contexts, even when texts dictated the way these objects were supposed to be used, they were not always followed in practice. Defining amulets solely as protective items limits our understanding of the diverse roles they performed in specific contexts where they were found. This nuanced approach is especially critical in Nubian studies, where moving beyond an Egyptocentric framework allows for a fuller recognition of the region’s internal diversity and distinct cultural practices. This is not to say that comparisons and Egyptian parallels are not relevant. The enduring entanglement of Egyptian and Nubia material culture, particularly during the high flow of foreign imports during the five-hundred-year New Kingdom conquest \parencites{lemosMaterialCultureColonization2020}, makes it counterproductive to ignore potential connections. As a result, the challenge is to understand how Egyptian-style amulets were adopted and adapted within the Nubian society. By drawing comparisons with contemporary Egyptian examples, we can better identify what is Nubian about the amulets’ role in local contexts.


Accordingly, this dissertation understands amulets as pendants designed to represent culturally recognizable motifs, such as deities and animals. As portable objects made from a range of varied materials, amulets were objects which could be acquired by different communities and, therefore, they performed different roles in each context.


Despite their prevalence and cultural significance, amulets have been the subject of few dedicated studies, mainly focused Egyptian contexts, particular motifs and museum collections \parencites{petrieAmulets1914, müller-winklerÄgyptischenObjektAmuletteMit1987, andrewsAmuletsAncientEgypt1994, herrmannÄgyptischeAmuletteAus1994, herrmannÄgyptischenAmuletteSammlungen2003, dubielAmuletteSiegelUnd2008, herrmann1001AmulettAltägyptischer2010, connorAmuletiDellAnticoEgitto2017, muñozpérezAllThatGlisters2018, muñozpérezAmulettesFunérairesÉgyptiennes2021, muñozpérezBúsquedaVidaEterna2021, hardwickLéventailCroyancesAmulettes2019, herrerínAnálisisPreliminarMomias2019, kalloniatisAmulets2019, lenzoFormulePourLamulette2019, arroyoPapelDosAmuletos2020, arroyoUsoAmuletosComo2021, donnatGestesRituelsAutour2020, donnatProposQuelquesAspects2021, mekisHypocephalusAncientEgyptian2020, quackBemerkungenAmulettpapyrusPBM2020, quackSprücheFürPerle2023, györyEmergenceBesAmulets2024}. A recent study by \textcite{quackAltägyptischeAmuletteUnd2022} investigates the chronological use of funerary amulets in Ancient Egypt, specifically examining instructions for their role in ritual texts, such as the Book of the Dead. Although not directly analyzing archaeological contexts and focusing on written sources only accessible to a small portion of the population, this study details state religion perspectives on the role of amulets.


Scholarly work, with the exception of excavation reports, has focused on specific periods and iconographic motifs \parencites{balandaNapatanFuneraryAmulets2015, balandaProtectingMummyReinterpretation2014, ditriaUnderstandingKermaAmulets2021, ditriaAmuletsKermaCulture2022, lemosForeignObjectsLocal2020, sackho-autissierQuelquesAmulettesNapatéennes2004}. A key contribution by \textcite[p.~39-52]{balandaDevelopmentRoyalFunerary2020} challenges “Egyptianization” idea by comparing Napatan amulet distribution in royal and non-royal burial sites. The author argues that non-royal amulet usage may have been rooted in local observations and narratives spread along the Nile Valley, and that the presence of Osirian amulets were a sign of that particular individual’s knowledge of religious texts. However, the study prioritizes religious beliefs as the primary motivating factor in the selection of amulets. This dissertation addresses this gap by investigating broader cultural and social dynamics that influenced amulet selection and deposition across different social groups, regions and phases within the Napatan period, integrating an analysis of the changing internal and external political contexts of the era.


\section{Research questions and objectives}


In this dissertation, I will address four key research questions through the analysis of the data:


\begin{enumerate}
	\item What were the kinds of Egyptian-style amulets used in Napatan Nubia in regards to their designs and materials?
	\item It is possible to identify local adoption and adaptation of certain Egyptian-style amulet motifs?
	\item What are the similarities and differences between amulet deposition in burials of royalty, elite and non-elite groups?
	\item Did amulet use change in these social groups depending on the chronological phase and geographically region?
\end{enumerate}


Based on these questions, the project tests the hypothesis that an objectscapes approach to the data can re-frame our understanding of Egyptian-style amulets in Napatan Nubia. Rather than viewing these objects only through a religious lens, we can consider them as goods which were manufactured, traded and consumed, and their meaning and use were actively shaped by the specific contexts of the social groups that had access to them. Therefore, by addressing these questions, the dissertation has the following objectives:


\begin{enumerate}
	\item Contribute to research into the Napatan period, by not only considering its influence over Upper Egypt, but also considering the internal diversity of Nubians and agency of its people.
	\item Examine how variability in the access to and use of amulets deposited as grave goods across different social groups in Napatan Nubia materializes structures of social complexity and reveals patterns of status differentiation.
	\item Shed light into the consumption and access to Egyptian-style amulets focusing on chronological changes, regional variation, and the role of Kushite rule over Egypt and geographic proximity of Nubian sites to the Egyptian border.
\end{enumerate}


Ultimately, the research has the aim to advance Nubian studies by moving from the focus of the external influence of Egypt to instead examine the internal social dynamics of Nubian communities. By analyzing the variability in the deposition of amulets in funerary contexts, this project investigates how differential access to goods materialized social and status differentiation within Napatan society. Furthermore, it explores how these consumption patterns were shaped by political shifts, regional variation, and the unique geopolitical context of Kushite rule over Egypt.


\section{Structure of the dissertation}


This dissertation is organized into six chapters. Following the introductory chapter, Chapter 2 centers on the theoretical and methodological frameworks employed throughout the research. In this chapter, I will discuss the shift in interpreting the presence of Egyptian-style material culture in Nubia from Egyptianization to cultural entanglements. Particularly, I will highlight the recent concept of objectscapes, which is particularly applicable to this research and considerations of foreign objects in local contexts, as well as their role in different social contexts.


Regarding the methodological aspects, further sections will present the organization of the database and the chronological division of the Napatan period employed in the thesis. Particularly, Chapter 2 also presents the reasons used to differentiate of social status based on burial customs. This discussion is crucial in order to distinguish the elite and non-elite tombs, in particular, which were located in the same cemeteries.


Chapter 3 focuses on the contextualization of the Napatan period, as well as the the history of cultural interactions between Nubia and Egypt from the 4th millennium BCE to the Napatan period. Additionally, this chapter will also explore questions regarding the origin and chronology of the Napatan dynasty, presenting the different interpretations posited by scholars.


The royal burial sites of El-Kurru and Nuri are the focus of Chapter 4, where  I analyze their burial customs chronologically, based on the division presented in Chapter 2. Furthermore, I will consider differences between different tomb owner’s royal title (king, queen and minor wife). The analysis are two fold: first, focusing on tomb structure and object types in order understand Napatan royal burial patterns; and second, exploring the distribution of amulets, their motifs and materials chronologically and by royal title.


Structurally, Chapter 5 parallels the previous chapter, but centers on the elite and non-elite sites. To establish a clear contrast, I will outline both elite and non-elite burial customs by their tomb structure and grave goods in order to distinguish them from each other and those of royalty. Afterwards, the amulets from both social groups will be analyzed chronologically and geographically, considering the same aspects employed in the analysis of royal amulets, i.e., distribution, motifs and material composition.


Chapter 6 applies the objectscapes approach to a comparative analysis of the amulets. It explores patterns of amulet consumption through three different criteria—social, chronological and regional. As in previous chapters, material characteristics are considered alongside contemporary parallels from Egypt, in order to identify local designs and adaptations.


Chapter 7 Discussion of the finds


Chapter 8 brings together the discussions and analysis presented throughout the dissertation, arguing that, in the context of Napatan Nubia, 
variations / internal diversity / local agency / adoption of Egyptian motifs and designs, but with local adaptations and novel designs / socially - / temporal - / regionally -